The overarching aim of the Local Langlands Correspondence (LLC) is to classify the smooth admissible complex representations of reductive $p$-adic groups in terms of geometric and number theoretic data. More concretely, if $G$ is a reductive $p$-adic group over some non-archimedean field $F$, the Langlands philosophy predicts that the irreducible (smooth admissible) complex representations $\Irr(G)$ should be controlled by the Weil group $\mathcal{W}_F$ of the field $F$, together with the geometry of the complex dual group $G^\vee$ of $G$. In \cite{Lusztig1995}, Lusztig defined the subset of irreducible unipotent representations $\Irr_{\un}(G)$. This is the smallest subset of $\Irr(G)$ containing all irreducible representations containing an Iwahori-fixed vector (also denoted Iwahori-spherical) and closed under the partition of $\Irr(G)$ into $L$-packets. More concretely, the Local Langlands Correspondence predicts the existence of a bijection,
\begin{align}
    \LLC_{\un}: \{(x,\phi)\ |\ x\in G^\vee, \phi\in\widehat{A_x}\}&\longleftrightarrow\Irr_{\un}(G),\\\notag
    (x,\phi)&\longmapsto\pi(x,\phi),
\end{align}
where $A_x$ is the component group of the centralizer $Z_{G^\vee}(x)$, and the left hand side can be seen to be in bijection with unramified Langlands parameters (trivial on the inertia subgroup). In \cite{KazhdanLusztig1987}, Kazhdan--Lusztig had already observed that some of the labels $\{(x,\phi)\ |\ x\in G^\vee, \phi\in\widehat{A_x}\}$ naturally parametrized Iwahori-spherical representations. The Langlands parametrization for the set of unipotent representations has come after many years of development. Lusztig's work \cite{Lusztig1995}, \cite{Lusztig2002} settled it for adjoint $p$-adic groups, Reeder \cite{Reeder2002} for Iwahori-spherical representations of split groups of arbitrary isogeny, and Solleveld for all reductive $p$-adic groups \cite{Solleveld2023},\cite{Solleveld2023b}. In addition, Lusztig characterized unipotent representations in terms of the structure of $G$ and representation theory of reductive quotients of its parahoric subgroups, which we revisit in Sections \ref{sec:finite_groups} and \ref{sec:padic_groups}. Deligne--Lusztig theory provides a powerful tool to study representations of finite groups of Lie type and, a posteriori, of $p$-adic groups.

While many questions in the unipotent LLC have now been resolved in full generality, some of them remain open. One such question relates to the stability of $L$-packets. Moeglin and Waldspurger tackled this problem for $G=\SO_{2n+1}(F)$ in \cite{MoglinWaldspurger2003} by considering the restriction of unipotent representations to maximal open compact subgroups. These subgroups provide a passage from unipotent representations of the $p$-adic group to unipotent representations of certain finite reductive groups. Work of Lusztig, Deligne and others in the late 20th century revealed remarkable structure on unipotent characters of these finite reductive groups, which Moeglin and Waldspurger exploited. Lusztig defined an involution on the space of virtual unipotent characters known as \textit{Lusztig's non-abelian transform} taking the irreducible characters to the so-called \textit{almost characters}. These characters can be obtained independently via the theory of character sheaves and, as a result, have nice geometrical and stability properties. 

The hope is that, if we can lift the Fourier transform to the setting of $p$-adic groups with analogous properties, one could use the resulting virtual characters to prove stability properties of $L$-packets. Moreover, it is believed by Lusztig that an analogous theory of character sheaves for $p$-adic groups should exist. Work of Ciubotaru \cite{Ciubotaru2020} and Aubert--Ciubotaru--Romano \cite{AubertCiubotaruRomano2024} has already considered this problem and formulated precise conjectures. We consider the same setting and similar notation to \cite{AubertCiubotaruRomano2024}, which we now describe in the simplified setting when $G$ is of simply connected type. 

Let $\max(G)$ denote the set of conjugacy classes of maximal open compact subgroups of $G$, which coincides with the set of conjugacy classes of maximal parahoric subgroups since $G$ is simply connected. For each $K\in\max(G)$, we define the \textit{parahoric restriction map} \textcolor{red}{Does it make much sense to avoid this in the introduction? Feels like I need to introduce $U_K$, $\overline{K}$ and the representation spaces anyway.}
\begin{align*}
    \res_{\un}^{K}:R_{\un}(G)&\longrightarrow R_{\un}(\overline{K})\\
    V&\longmapsto V^{U_K},
\end{align*}
where $U_K$ is the pro-$p$ unipotent radical $K$, $\overline{K}:=K/U_K$ is a connected finite group of Lie type and $R_{\un}(G)$ and $R_{\un}(\overline{K})$ are the complexification of the Grothendieck groups spanned by the unipotent representations of $G$ and $\overline{K}$, respectively. Secondly, Lusztig's nonabelian Fourier transforms on $R_{\un}(\overline{K})$ for each $K\in\max(G)$ gives an involution 
\begin{equation*}
    \FT^{\cpt}_{\un}:\bigoplus_{K\in\max(G)}R_{\un}(\overline{K})\longrightarrow\bigoplus_{K\in\max(G)}R_{\un}(\overline{K}).
\end{equation*}

The remaining ingredient to formally state Aubert--Ciubotaru--Romano's conjecture \cite{AubertCiubotaruRomano2024}*{Conjecture 1.2} is the construction of the \textit{dual nonabelian Fourier transform} $\FT^\vee:R_{\un}(G)\longrightarrow R_{\un}(G)$ compatible with the unipotent restriction maps $\res_{\un}^K$ and Lusztig's nonabelian Fourier transform $\FT^{\overline{K}}$. The existence of such a map is conjectured, but so far no explicit general construction has been achieved on the infinite-dimensional space $R_{\un}(G)$. Instead, Ciubotaru--Opdam \cite{CiubotaruOpdam2010} realized that, for applications, it was often sufficient to study the quotient space $\overline{R}_{\un}(G)$ of $R_{\un}(G)$ by the span of all properly parabolically induced representations, since parabolic induction is generally well understood. The upshot of this simplification is double: firstly, the quotient space $\overline{R}_{\un}(G)$ is finite dimensional, in contrast to the infinite-dimensional space $R_{\un}(G)$. Secondly, $\overline{R}_{\un}(G)$ is canonically isomorphic (in fact, isometric with respect to some canonical inner products) to the elliptic unipotent representation space $R_{\un,\mathrm{ell}}(G)$, a certain subspace of $R_{\un}(G)$. Moreover, the Local Langlands Correspondence provides an explicit description for a natural orthogonal basis of $R_{\un,\mathrm{ell}}(G)$. To describe this basis, let $u\in G^\vee$ be a unipotent element and let $\Gamma_u$ be the reductive part of the centralizer $Z_{G^\vee}(u)$. We then consider the set of elliptic pairs 
$$\mathcal{Y}(\Gamma_u)_{\mathrm{ell}}=\{(s,h)\in\Gamma_u\ |\ sh=hs\text{ and }Z_{\Gamma_u}(s,h)\text{ finite}\}$$
and, for each such elliptic pair $(s,h)$, we define the virtual characters 
$$\Pi(u,s,h):=\sum_{\phi\in A_{us}}\phi(h)\pi(us,\phi).$$
The Local Langlands Correspondence induces an isomorphism
\begin{equation*}
    \mathrm{LLC}^{\mathrm{ell}}_{\un}:\bigoplus_{u\in G^\vee}\CC[\mathcal{Y}(\Gamma_u)]^{\Gamma_u}\longrightarrow R_{\un,\mathrm{ell}}(G),\quad\quad (u,s)\longmapsto\Pi(u,s,h),
\end{equation*}
where the sum in the left hand side is taken over conjugacy classes of unipotent elements in $G^\vee$, so in particular $R_{\un,\mathrm{ell}}(G)$ is the span of $\{\Pi(u,s,h)\ | \ u\in G^\vee\text{ unipotent, }(s,h)\in\mathcal{Y}(\Gamma_u)_{\mathrm{ell}}\}$, where the parameters are taken up to $G^\vee$-conjugation.

The space of elliptic pairs has an obvious involution given by the flip $(s,h)\mapsto(h,s)$, and this defines an involution 
$$\FT^\vee_{\mathrm{ell}}:R_{\un,\mathrm{ell}}(G)\longrightarrow R_{\un,\mathrm{ell}}(G),$$
the dual elliptic nonabelian Fourier transform. We are finally ready to state the main conjecture of this paper, in the simplified setting of split, simply connected $p$-adic groups.
\begin{conjecture}\label{conj:main}\cite{AubertCiubotaruRomano2024}*{Conjecture 1.2}
    Let $G$ be a split semisimple simply connected $p$-adic group. Then the diagram 
    \[
        \begin{tikzcd}
        R_{\un,\text{ell}}(G) \arrow[r, "\FT^\vee_{\text{ell}}"] \arrow[d, "\res_{\un}^{\cpt}"] & R_{\un,\text{ell}}(G) \arrow[d, "\res_{\un}^{\cpt}"] \\
        \bigoplus_{K\in\max(G)}R_{\un}(\overline{K}) \arrow[r, "\FT^{\cpt}_{\un}"] & \bigoplus_{K\in\max(G)}R_{\un}(\overline{K}),
        \end{tikzcd}
    \]
    commutes, up to certain roots of unity.
\end{conjecture}

In \cite{AubertCiubotaruRomano2024}, Aubert--Ciubotaru--Romano stated a more general conjecture including all isogeny classes, but the technical details are unnecessary for our purposes. Moreover, they showed that this conjecture (or the more general version) holds for $G=\SL_n(F)$,  $\PGL_n(F)$ or $\Sp_4(F)$. In previous work, Ciubotaru and Opdam \cite{Ciubotaru2020} showed it holds for $G_2(F)$, and Waldspurger \cite{Waldspurger2018} showed it for $G=\SO_{2n+1}(F)$. In this document we settle the above conjecture for two further exceptional cases. 


\begin{theorem} Let $G$ be a simple, split, simply connected $p$-adic group of type $F_4$ or $E_6$. Then Conjecture \ref{conj:main} holds for $G$.
\end{theorem}

\subsection*{Structure of the document}
The document is organized as follows. We devote Section \ref{sec:finite_groups} to finite groups of Lie type and Deligne--Lusztig theory, with a focus on the construction of Lusztig's nonabelian Fourier transform. In Section \ref{sec:padic_groups} we discuss representation theoretic aspects of $p$-adic groups and their relation to representations of finite groups of Lie type. We recall the definition of unipotent representations and construct the dual nonabelian Fourier transform. Section \ref{sec:parahoric_restriction} is the most technical; there we discuss in detail how to explicitly compute the restriction functors $\res_{\un}^{\cpt}$, a method developed by Reeder in \cite{Reeder2000} for adjoint groups and which we adapt for simply connected groups. Sections \ref{sec:F4} and \ref{sec:E6} contain geometric information about the complex reductive groups $G^\vee=F_4(\CC)$ or $E_6(\CC)_{\scn}$ necessary to explicitly describe the virtual representations $\Pi(u,s,h)$ spanning the elliptic unipotent representation space $R_{\un,\mathrm{\ell}}(G)$. 

In Section \ref{sec:tables}, we display all tables containing the parahoric restrictions of unipotent representations with respect to maximal parahoric subgroups. From these tables, the proof of Conjecture \ref{conj:main} is a simple verification.


